\section{从狄道到建业：征服怎样变成再统一}

\subsection{261--265：蜀亡与魏易}

二六二年姜维\eventanchor{TK-0262-JIANG-WEI-DIDAO}{围狄道，邓艾据险救援，
蜀军撤退}。这次失利不能单独解释蜀亡，却再次显出北伐的基本代价：陇右或许
给小国以改变边境的机会，粮运和魏军的纵深却会把机会转为消耗。

二六三年司马昭发动\eventanchor{TK-0263-WEI-INVADES-SHU}{多路伐蜀}，钟会、
邓艾、诸葛绪由关中、陇右南下。剑阁方向的牵制和阴平方向的突入共同压缩了
成都的选择；邓艾抵成都平原后，
\eventanchor{TK-0263-CHENGDU-SURRENDER}{刘禅出降}。可确认的结果不是“蜀人
一夜之间放弃抵抗”，而是都城缺少可及时会合的野战军，援军也难以穿过已经被
切断或牵制的走廊。阴平行军的细节、人数和城中劝降的每一句话，材料并不允许
写成确定的现场记录。

灭蜀同时改变了魏的名义结构：曹奂在十二月
\eventanchor{TK-0263-SIMA-ZHAO-DUKE-JIN}{进封司马昭为晋公、加九锡}。次年
司马昭又\eventanchor{TK-0264-SIMA-ZHAO-KING-JIN}{进为晋王}；成都却发生
\eventanchor{TK-0264-ZHONG-HUI-REBELLION}{钟会之乱}，邓艾被杀，钟会亦为
魏军与降众所杀。征服者可以取得降表，却未必立刻掌握军队、降众和地方消息；
这场兵变正是占领与接收之间的裂缝。

江东也在二六四年更换皇帝，
\eventanchor{TK-0264-SUN-HAO-ACCESSION}{孙皓继立}。一年后
\eventanchor{TK-0265-SIMA-ZHAO-DEATH}{司马昭去世，司马炎继承晋王位和实权}，
曹奂随即\eventanchor{TK-0265-WEI-ABDICATION}{禅位，西晋建立}。魏晋禅代与
汉魏禅代一样，能够确定礼仪的顺序；它不能仅凭诏文告诉我们所有参与者如何
判断风险。两次交接都依赖此前已在中枢积累的军政控制。

\subsection{268--274：准备渡江，而非等待它自然发生}

西晋二六八年\eventanchor{TK-0268-TAISHI-CODE}{颁行《泰始律》}，以十八篇
新律整理魏晋以来的刑法框架。律文说明统一中的新朝希望提供可援引的规则，
并不证明州郡已经同速执行。次年羊祜
\eventanchor{TK-0269-YANG-HU-JINGZHOU}{都督荆州，在襄阳经营屯戍、招纳和
对吴关系}；这是把未来战事化为长期边境经营，而不是把长江看成一条可在某日
突然越过的线。

二七二年西陵督步阐降晋，陆抗
\eventanchor{TK-0272-XILING}{围西陵、击退晋援，收复峡口}。吴仍能调动一套
有韧性的江防网络，西晋也因此看见单纯从荆州推进的限度。二七四年
\eventanchor{TK-0274-LU-KANG-DEATH}{陆抗去世}，吴失去熟悉西部边境的统帅；
人事损失不是自动的崩溃，却使边防、朝廷和继承问题少了一处可以调节的接口。

\subsection{278--280：多路进军与一份降表之后}

羊祜在二七八年临终前上表建议伐吴，
\eventanchor{TK-0278-YANG-HU-DEATH}{杜预接掌荆州部署}。到二七九年十一月，
\eventanchor{TK-0279-JIN-INVADES-WU}{杜预、王濬、王浑等分路伐吴}：益州水军
沿江而下，荆、淮军并进。它之所以能够发生，靠的是此前对蜀地、荆州、船队和
边境信息的多年接收；它仍不是一支舰队单独完成的故事。

二八〇年王濬舰队抵建业，
\eventanchor{TK-0280-SUN-HAO-SURRENDER}{孙皓出降}，孙吴灭亡。全国统一在
政权名义上完成，但共同年代的终点不是一张从此没有摩擦的地图。新的朝廷还要
接收江东的官吏、道路、簿籍、港口和地方关系；这些不同速度的接收，正是下文
回看三国国家、跨境走廊与社会材料的理由。

\begin{turningpoint}[二八〇年的终点]
灭吴的直接成果是最高权属的合并；中期问题则是西晋如何把蜀、吴旧有的行政和
地方网络接入自己的任命、法律和赋役语言。把“统一”同时读作军事、制度和
日常生活的不同速度，才能避免把降表当成一切已经完成。
\end{turningpoint}

\begin{sourcenote}[征服、禅代与法律的证据]
灭蜀、魏晋禅代和灭吴主要依《三国志》《晋书》与《资治通鉴》互证；官修帝纪
和禅诏最适于确认顺序与名义，不足以单独还原协商。关于《泰始律》应以《晋书》
刑法志及法律史线索讨论其编成，不能把法典文本等同于全国执行。
\end{sourcenote}
